	
\section*{Abstract}
		This paper provides a detailed analysis of the methodological principles of theoretical physicist John F. Donoghue, as expressed in his recent interview \cite{DonoghueInterview2025} and his publications. It compares his willingness to question established assumptions—the claimed fundamental incompatibility of quantum mechanics and gravity, the principle of Naturalness, and the unification bias—with the approach of the Fundamental Fractal-Geometric Field Theory (FFGFT) within the T0 theory. The study sets out that Donoghue's approaches to the effective field theory (EFT) of gravity, quadratic gravity, and random dynamics are methodologically compatible with the revisions of the FFGFT. Donoghue himself does not comment on the FFGFT; the parallels are an assessment made in this document. It shows how the FFGFT builds a dynamic vacuum field $\Phi(x)$ from the T0 time-mass duality and a fractal geometry, and thereby undertakes a re-evaluation of fundamental assumptions of the kind Donoghue considers legitimate. A key focus is that the FFGFT framework is built up step by step from simple core structures (Dirac form, Lagrangian), i.e. it follows a bottom-up principle that fits Donoghue's skepticism toward top-down unification.

		\medskip\noindent\textbf{Status markers.} Statements marked \textbf{[S]} are \emph{speculative}: a hint or conjecture that has not yet been derived or numerically checked in the corpus. The remaining statements of this document are methodological assessments and comparisons; it contains no new computational results.

	
	
	
	Modern theoretical physics stands at a crossroads. While the Standard Model of particle physics and General Relativity (GR) make exceptionally precise predictions within their respective domains, fundamental questions regarding their unification, the quantization of gravity, and the nature of space, time, and vacuum remain unanswered. Within this landscape of potential solutions, John F. Donoghue occupies a clearly defined position. His work is characterized less by new postulates than by a consistent \emph{methodological revisionism}: the systematic questioning and, where necessary, abandonment of assumptions that prove to be obstacles to consistent theoretical progress.
	
	The Fundamental Fractal-Geometric Field Theory (FFGFT), embedded within the framework of the T0 theory, proposes a different path. Rather than treating gravity as an irreducible geometric property of spacetime, it models it as an emergent phenomenon arising from perturbations of a fundamental, dynamic vacuum field $\Phi(x) = \rho(x) e^{i\theta(x)}$, whose structure is determined by an underlying fractal geometry. This approach requires the explicit revision of several central pillars of modern physics: the concept of a passive vacuum, the notion of gravity as primary geometry, and the expectation of a top-down unification through extended symmetries.
	
	This paper examines how far Donoghue's principles, as formulated particularly in his interview \cite{DonoghueInterview2025}, agree with the methodological approach of the FFGFT. We first present the core arguments of the FFGFT, then Donoghue's positions with special consideration of the interview statements, and finally the methodological and content-related parallels. These parallels are our own assessment; they are not a statement by Donoghue on the FFGFT.
	
	
	The FFGFT builds on the axioms of the T0 theory, the core of which is the time-mass duality
	\[
	T(x,t) \cdot m(x,t) = 1.
	\]
	This duality establishes an intrinsic, reciprocal relationship between temporal and massive degrees of freedom and opens an approach to gravity, understanding it as an effect of vacuum dynamics within a fractal background geometry. The fundamental dimensionless constant $\xi$ of the T0 theory is interpreted here as the \emph{intrinsic fractal packing deficit of three-dimensional Euclidean space}, which gives the theory its name.
	
\section{The Dynamic Vacuum Field \texorpdfstring{$\Phi(x)$}{Phi(x)} as Fundamental Substance}
	
	The central object is the complex scalar field
	\[
	\Phi(x) = \rho(x) e^{i\theta(x)},
	\]
	which does not represent a particle field within the vacuum, but the physical vacuum \emph{itself}. Its components have a clear phenomenological interpretation:
	\begin{itemize}
		\item $\rho(x)$: The vacuum amplitude, directly correlated with the massive component of the T0 duality: $m(x,t) = 1/T(x,t)$.
		\item $\theta(x)$: The vacuum phase, whose dynamics follow from the rotation of T0 node structures and endow the vacuum with an intrinsic temporal rhythm.
	\end{itemize}
	
	The undisturbed ground state is $\Phi_0 = \rho_0 e^{-i\mu t}$, with the fundamental scale $\rho_0 = 1/\xi^2$ fixed by the T0 geometry and the frequency $\mu = \xi m_0$. This gives the vacuum a natural ``pacemaker'' with $\dot{\theta} = m = 1/T$.
	
	\section{Mathematical Core: Derivation from Simplified Dirac and Lagrangian Structures}
	
	The FFGFT does not postulate its final, complex field-theoretic framework axiomatically. Instead, it is built up step by step (\emph{bottom-up}) from the simplest mathematical structures of the T0 core. The following steps are to be read as a construction path; where an assignment is set rather than derived, this is stated.
	
	\begin{enumerate}
		\item \textbf{Starting Point - T0 Core Axioms:}
		\begin{itemize}
			\item Fundamental time-mass duality: \( T(x,t) \cdot m(x,t) = 1 \).
			\item Associated simplified geometric constant \( \xi \), interpreted as an intrinsic fractal packing parameter.
		\end{itemize}
		
		\item \textbf{First-Level Derivation - Simplified Quantum Dynamics:}
		\begin{itemize}
			\item From the duality, a \textbf{simplified form of the Dirac equation} is derived. This step connects the classical duality to a quantum-mechanical operator structure, establishing a bridge to quantum field theory without initially requiring its full complexity.
			\item A \textbf{simplified Lagrangian}, e.g., \( \mathcal{L}_{\text{simple}} \propto (\partial \Delta m)^2 \), is constructed. It describes the dynamics of deviations (\( \Delta m = m - m_0 \)) from the vacuum mass configuration \( m_0 \).
		\end{itemize}
		
		\item \textbf{Second-Level Derivation - Emergence of the Full Field Theory:}
		\begin{itemize}
			\item The degrees of freedom of the simplified framework are mapped to the components of the \textbf{complex vacuum field} \( \Phi(x) = \rho(x) e^{i\theta(x)} \) (this assignment is a setting):
			\begin{align*}
				\rho(x) &\leftrightarrow m(x,t) = 1/T(x,t) \quad \text{(amplitude from mass density)} \\
				\theta(x) &\leftrightarrow \text{phase from rotational dynamics of T0 ``nodes''}.
			\end{align*}
			\item With this mapping, the \textbf{FFGFT Lagrangian} follows:
			\[
			\mathcal{L}_{\text{FFGFT}} = (\partial \rho)^2 + \rho^2 (\partial \theta)^2 - \frac{1}{2} m_T^2 (\rho - \rho_0)^2 + \xi m_\ell \bar{\psi}_\ell \psi_\ell \Delta m + \dots
			\]
			The kinetic term \( (\partial \rho)^2 + \rho^2 (\partial \theta)^2 \) is the direct image of the simplified term \( (\partial \Delta m)^2 \) within the new field formalism.
		\end{itemize}
	\end{enumerate}
	
	Along this path, the field-theoretic level of description (the FFGFT) is not an independent postulate but is tied to the simpler T0 foundations; the mapping in step~3 remains a setting.
	
	\section{Lagrangian Formulation from T0 Principles}
	
	The Lagrangian of the FFGFT, as it results from the construction above, is:
	
	\begin{align}
		\mathcal{L}_{\text{FFGFT}} &= \underbrace{(\partial_\mu \rho)(\partial^\mu \rho) + \rho^2 (\partial_\mu \theta)(\partial^\mu \theta)}_{\text{Kinetic Terms from T0 Mapping}} \\
		&\quad - \underbrace{\frac{1}{2} m_T^2 (\rho - \rho_0)^2}_{\text{Potential from T0 Mediator Mass } (m_T = \lambda / \xi)} \\
		&\quad + \underbrace{\xi m_\ell \bar{\psi}_\ell \psi_\ell \Delta m}_{\text{Matter-Vacuum Coupling}} + \cdots
	\end{align}
	
	Here, $\Delta m = m - m_0$ denotes the deviation from the vacuum mass configuration. This Lagrangian describes how matter ($\psi_\ell$) locally perturbs the vacuum field $\Phi$ and how these perturbations propagate and interact.
	
	\section{Approaches to Open Problems}
	
	The fractal-geometric framework of the FFGFT yields approaches to three open problems. They are worked out to different degrees:
	
	\begin{enumerate}
		\item \textbf{Gravity as Vacuum Convergence}: Instead of abstract spacetime curvature, gravity arises through the local convergence and densification of the vacuum field $\Phi$ in response to material stress-energy. The observed geometry is emergent.
		\item \textbf{Singularity-Free Black Holes}: Black holes appear as stable, highly condensed configurations of T0 nodes in the vacuum field. The GR singularity is an artifact of the classical, effective description that neglects the underlying regular fractal T0 structure.
		\item \textbf{Cosmology without Inflation and Dark Energy} \textbf{[S]}: The infinite homogeneous T0 geometry with its fractal scale $\xi_{\text{eff}} = \xi/2$ could provide an alternative mechanism for the observed cosmic acceleration and the CMB anisotropies, without inflation fields or a cosmological constant. This point is speculative; in the current corpus the cosmic sector is deliberately left aside, because the Standard Model does not derive it either (P39).
	\end{enumerate}
	
	
	Donoghue's contributions to theoretical physics are characterized less by a comprehensive theory of his own than by a clear methodological stance. His positions, as evident in his interview \cite{DonoghueInterview2025} and his writings \cite{Donoghue1995, Donoghue2022}, can be summarized in four core principles.
	
	\section{Principle 1: Effective Field Theory as a Universal and Sufficient Framework}
	
	Donoghue views both GR and the Standard Model unequivocally as \emph{effective field theories} (EFTs)—theories that are valid only up to a certain energy scale, beyond which new physics and new degrees of freedom become relevant.
	
	In the interview he states this unequivocally: \emph{``I think the popular phrasing is totally wrong, that quantum physics and gravity go perfectly well, as well as any other theory that we know about. Quantum gravity involves a field, which is the metric. That field is quantized. It was done by Feynman and DeWitt in exactly the same way we do QCD; there's no difference at all in the framing of it.''} \cite{DonoghueInterview2025} (04:31-05:10).
	
	This position dismantles the widespread narrative of a fundamental incompatibility. The supposed problems of quantum gravity—particularly non-renormalizability—are, according to Donoghue, not fatal flaws but natural \emph{hints} at the limits of GR as an EFT and the need for new physics at the Planck scale \cite{Donoghue1995}. The EFT framework allows for precise quantum field theoretical predictions to be made within the known theory (as his calculation of quantum corrections to the Newtonian potential demonstrates), without knowledge of the ultimate UV completion.
	
	\section{Principle 2: Pragmatic Renormalizability through Axiom Revision (Quadratic Gravity)}
	
	As a minimalist and "conservative" extension of GR, Donoghue advocates for \emph{quadratic gravity}, where terms like $R^2$ and $R_{\mu\nu}R^{\mu\nu}$ are added to the Einstein-Hilbert action \cite{Salvio2018}. This theory is renormalizable, as shown by Stelle in the 1970s, but requires relinquishing the established principle of microcausality at high energies.
	
	In the interview he explains this compromise: \emph{``The nature of the theories with higher derivatives is that you get a massless [...] particle with the usual arrow of causality and a very heavy particle with the opposite arrow of causality.''} \cite{DonoghueInterview2025} (34:16-36:45). These "dueling arrows of causality"—the existence of a ghost degree of freedom that effectively propagates backward in time—is accepted by Donoghue as a legitimate price for a mathematically consistent (renormalizable) quantum theory of gravity. This stance shows a clear prioritization of \emph{mathematical consistency} and \emph{empirical adequacy} (the theory is identical to GR at low energies) over the strict adherence to all traditional axiomatic requirements.
	
	\section{Principle 3: Skepticism towards ``Naturalness'' and the Unification Bias}
	
	Donoghue subjects two guiding principles of particle physics to fundamental criticism: the principle of Naturalness and the belief in a Grand Unified Theory (GUT).
	
	He argues that the failure to find supersymmetry at the LHC dealt a severe blow to the Naturalness argument that drove the search for new physics for decades \cite{DonoghueInterview2025} (47:51-50:04). More fundamentally, he criticizes the \emph{unification bias}: \emph{``We've never really seen unification. [...] The idea of unification could just totally be a bias.''} \cite{DonoghueInterview2025} (44:22-45:12). He sharply distinguishes between the successful \emph{merger} of seemingly different phenomena (like electricity and magnetism) under a common theoretical structure and the speculative \emph{unification} of separate interactions (strong, weak, electromagnetic) into a single larger symmetry group, for which there is no empirical evidence.
	
	\section{Principle 4: ``Random Dynamics'' and Anti-Unification as an Alternative Paradigm}
	
	As a conceptual counter-proposal, Donoghue favors Holger Nielsen's idea of \emph{Random Dynamics} \cite{DonoghueInterview2025} (41:57-43:38). This scenario posits that at extremely high energies, initially "everything possible" exists. Only certain structures—those "protected" by symmetries like gauge invariance, chirality, and general covariance—are robust enough to survive down to the low-energy scales we observe.
	
	This is the exact opposite of a traditional unification program. It is an \emph{anti-unification} or a \emph{bottom-up selection principle}: Instead of descending from an elegant, unified high-energy theory, one starts with a chaotic high-energy ``swamp'' and observes which structures endure into the low-energy regime through selective stability. Donoghue values this approach because it exemplarily shows how deeply rooted theoretical preferences (for elegance and symmetry) could distort our expectations of the fundamental theory.
	
	
	There are methodological parallels on several levels between Donoghue's revisionist approach and the basic conception of the FFGFT. The following table summarizes these parallels systematically.
	
	% ==============================================================================
	% LONGTABLE - English (Extended)
	% ==============================================================================
	\begin{longtable}{p{0.17\textwidth} p{0.38\textwidth} p{0.38\textwidth}}
		\caption{Systematic Comparison of the Methodological Principles of John F. Donoghue with their Correspondence and Application in the Fundamental Fractal-Geometric Field Theory (FFGFT)} \\
		\toprule
		\textbf{Conceptual Level} & \textbf{Donoghue's Principle and Argumentation} & \textbf{Correspondence and Application in T0/FFGFT} \\
		\midrule
		\endfirsthead
		
		\multicolumn{3}{c}{{\tablename\ \thetable{} -- continued}} \\
		\toprule
		\textbf{Conceptual Level} & \textbf{Donoghue's Principle and Argumentation} & \textbf{Correspondence and Application in T0/FFGFT} \\
		\midrule
		\endhead
		
		\midrule \multicolumn{3}{r}{{continued on next page}} \\
		\endfoot
		
		\bottomrule
		\endlastfoot
		% ==============================================================================
		% Table Content
		% ==============================================================================
		\textbf{1. Theory Limits and Revisions} & 
		\textbf{EFT Perspective}: GR and SM are effective theories with inherent limits. Their form (e.g., non-renormalizability) points towards new physics. He considers the widespread thesis of incompatibility false. \cite{DonoghueInterview2025} (04:31-05:18, 06:48-07:30) & 
		GR and QFT appear as \emph{low-energy effective limits} of T0 dynamics. The FFGFT proposes the dynamic fractal-geometric vacuum field $\Phi$ as the \emph{``new physics''} beyond the Planck scale. Abandoning the passive vacuum is the corresponding revision within this framework. \\
		\midrule
		
		\textbf{2. Prioritizing Mathematical Consistency} & 
		\textbf{Quadratic Gravity}: Renormalizability can be a higher good than strict adherence to microcausality at high energies. Pragmatic compromise in favor of a consistent quantum theory. \cite{DonoghueInterview2025} (34:16-36:45) & 
		Building a self-contained field theory from T0 principles prioritizes \emph{internal mathematical and conceptual consistency} of the whole system. The revision of established axioms (passive vacuum, geometry as cause) is accepted for the sake of this consistent overall picture. \\
		\midrule
		
		\textbf{3. Criticism of Established Dogmas} & 
		\textbf{Naturalness \& Unification Bias}: He regards Naturalness as largely a human prejudice (LHC findings) and the expectation of a Grand Unified Theory (GUT) as a bias without empirical basis. \cite{DonoghueInterview2025} (44:22-50:04) & 
		T0/FFGFT rejects \emph{Naturalness as a guiding principle}. Fine-tunings arise from the underlying universal fractal geometry ($\xi$). Unification is sought not through abstract symmetries (SUSY/GUTs) but through a common dynamic substrate ($\Phi$) from which the phenomena are to be derived. \\
		\midrule
		
		\textbf{4. Alternative Unification Path} & 
		\textbf{Random Dynamics / Anti-Unification}: Low-energy physics (SM+GR) as a robust, symmetry-protected remnant of an original high-energy random dynamics. \cite{DonoghueInterview2025} (41:57-43:38) & 
		Unification in T0/FFGFT follows a \emph{``Bottom-up'' principle}: From one basic axiom (time-mass duality) and a fractal base geometry, a field theory (FFGFT) is built up. This is a structural analogue to the selection in Random Dynamics. \\
		\midrule
		
		\textbf{5. Bottom-Up Construction} & 
		\textbf{Random Dynamics / Emergence}: Low-energy physics emerges as a stable structure from a simpler or chaotic high-energy starting point. Complexity is built, not assumed. \cite{DonoghueInterview2025} (41:57-43:38) & 
		\textbf{Construction from Simplified T0 Core}: The FFGFT (field $\Phi$, Lagrangian) is built up step by step from minimal axioms (time-mass duality) via simplified structures (Dirac form, simple Lagrangian). Unification is the result, not the starting point. \\
		
	\end{longtable}
	
	\section{Conceptual Parallels}
	
	\subsection{Rethinking Gravity as a Field Theory}
	
	Donoghue stresses that quantum gravity is a field theory like any other. This fits the core of the FFGFT: if GR geometry is a low-energy effective description, it is natural to look for an underlying field-theoretic microstructure. The FFGFT proposes the vacuum field $\Phi$ for this, whose perturbations and convergences are to generate the observed curvature. Donoghue's work shows that such a field-theoretic approach to gravity is methodologically not excluded; whether the FFGFT implementation holds up is decided by its own results.
	
	\subsection{Singularities as Artifacts of Effective Descriptions}
	
	Donoghue's EFT perspective provides a clear explanation for why GR singularities need not pose an insurmountable fundamental problem: GR is a \emph{low-energy effective theory} that exceeds its validity limit at the extreme densities in the center of a black hole. The FFGFT takes up this insight and models black holes as \emph{stable, singularity-free configurations} of T0 nodes in the vacuum field. In this reading, the singularity is an artifact of the incomplete effective description (GR).
	
	\subsection{Bottom-Up Derivation as Operationalization of Donoghue's Principles}
	
	The construction path of the FFGFT corresponds to the methodological preferences expressed by John F. Donoghue, particularly his skepticism toward top-down unification.
	
	\subsubsection{From Simplified Core to Emergent Complexity}
	Donoghue's affinity for "Random Dynamics" favors scenarios where the observed low-energy structure (like the Standard Model) is a robust remnant emerging from a simpler or even chaotic high-energy starting point \cite{DonoghueInterview2025} (41:57-43:38). The T0/FFGFT framework follows the same "bottom-up" logic:
	\begin{itemize}
		\item \textbf{Start Simple}: The theory begins with the minimal set of axioms (time-mass duality, fractal geometry \(\xi\)).
		\item \textbf{Build Up, Don't Postulate}: The field-theoretic apparatus (the complex field \(\Phi\), its Lagrangian, and its coupling to matter) is built up step by step from that simple core via intermediate simplified structures (Dirac form, simple Lagrangian); the field assignment is a setting.
		\item \textbf{Emergent Unification}: The unification of phenomena (gravity as vacuum convergence) is therefore not the starting assumption but the \emph{final result} of this derivation. This stands in contrast to top-down unification programs that begin with a large symmetry and attempt to derive the low-energy world from it.
	\end{itemize}
	
	\subsubsection{Consistency through Derivation}
	Donoghue prioritizes pragmatic mathematical consistency. In the FFGFT, this consistency is not meant to be enforced ad hoc but to follow from the construction path: the levels from the simplified T0 core to the field-theoretic formulation are intended as \emph{the same theory} in different mathematical languages. Whether they actually fit together without contradiction has to be checked at the individual transitions; the corpus verification scripts do this for the quantities treated in each case.
	
	\subsection{A Bottom-Up Path to Unification}
	
	Donoghue's sympathy for Random Dynamics and his criticism of the GUT bias make the alternative unification path of the T0 theory methodologically comprehensible. Instead of unifying all forces through ever-larger symmetry groups (as in SUSY or String Theory)—a top-down approach—the FFGFT attempts to derive the physical phenomena from a common \emph{geometro-dynamic substrate} (the vacuum field $\Phi$), whose properties are fixed by the T0 duality and the fractal geometry. This corresponds to the bottom-up or selection principle of Random Dynamics: from a simple, fundamental initial state, the complex structures of observed physics develop through internal dynamics.
	
	\section{Concrete Applications: Donoghue's Principles in FFGFT Argumentation}
	
	\subsection{Placing the Vacuum Field Revision}
	Donoghue's pragmatic stance in the Quadratic Gravity debate (giving up microcausality for renormalizability) shows that revising a principle considered fundamental is a recognized theoretical tool. This fits the central revision of the FFGFT: abandoning the passive vacuum concept of QFT in favor of an active, dynamic field $\Phi$ that represents the actual physical substance. Both revisions follow the same logic: they sacrifice a traditional axiom to achieve a higher theoretical goal (renormalizability or a unified description from time-mass duality).
	
	\subsection{Empiricism vs. Speculative Elegance}
	Donoghue's criticism of Naturalness and his own field after the LHC is a call for stricter empiricism. The FFGFT does not start from aesthetic or ``natural'' extensions of the Standard Model (like SUSY), but from a minimal principle (time-mass duality), and builds from it a concrete, calculable theory whose results are tested against measured values. The focus lies on internal consistency and the derivation of phenomena, not on fulfilling external notions of elegance.
		\begin{thebibliography}{99}
	
	% Primary source: Video interview
	\bibitem{DonoghueInterview2025}
	J. F. Donoghue,
	Interview: \emph{The Physicist Who Says We've Already Quantized Gravity},
	\textit{Theories of Everything with Curt Jaimungal}, YouTube, 2025.
	[Online]. Available: \href{https://www.youtube.com/watch?v=dG_uKJx6Lpg}{https://www.youtube.com/watch?v=dG\_uKJx6Lpg}.
	
	% Fundamental works by Donoghue on EFT of gravity
	\bibitem{Donoghue1995}
	J. F. Donoghue,
	\emph{Introduction to the Effective Field Theory Description of Gravity},
	in Advanced School on Effective Theories, Almunecar, Spain, 1995,
	arXiv:\href{https://arxiv.org/abs/gr-qc/9512024}{gr-qc/9512024}.
	
	\bibitem{Donoghue2022}
	J. F. Donoghue,
	\emph{Quantum General Relativity and Effective Field Theory},
	arXiv:\href{https://arxiv.org/abs/2211.09902}{2211.09902v1 [gr-qc]}, Nov. 2022.
	
	% Quadratic Gravity
	\bibitem{Salvio2018}
	A. Salvio,
	\emph{Quadratic Gravity},
	Front. in Phys., vol. 6, p. 77, 2018.
	doi: \href{https://doi.org/10.3389/fphy.2018.00077}{10.3389/fphy.2018.00077}.
	
	\bibitem{QuantaQuadratic2025}
	\emph{Old 'Ghost' Theory of Quantum Gravity Makes a Comeback},
	Quanta Magazine, Nov. 2025.
	[Online]. Available: \href{https://www.quantamagazine.org/old-ghost-theory-of-quantum-gravity-makes-a-comeback-20251117/}{https://www.quantamagazine.org/old-ghost-theory-of-quantum-gravity-makes-a-comeback-20251117/}.
	
	% Conceptual background: Randomness and Fractals
	\bibitem{GambiniPullin2005}
	R. Gambini and J. Pullin,
	\emph{Fundamental gravitational limitations to quantum computing},
	arXiv:\href{https://arxiv.org/abs/quant-ph/0507262}{quant-ph/0507262v1}, Jul. 2005.
	
	\bibitem{Bourgade2022}
	P. Bourgade and X. Chen,
	\emph{Liouville quantum gravity from random matrix dynamics},
	arXiv:\href{https://arxiv.org/abs/2206.03029}{2206.03029v2 [math.PR]}, Jul. 2025.
	
	% Standard reference literature
	\bibitem{WikiQG}
	Wikipedia, \emph{Quantum Gravity}.
	[Online]. Available: \href{https://en.wikipedia.org/wiki/Quantum_gravity}{https://en.wikipedia.org/wiki/Quantum\_gravity}.
	
	% Primary source of the T0/FFGFT theory
	\bibitem{PascherT0Book}
	J. Pascher,
	\emph{T0 – Time-Mass Duality Part 1: Foundations – From Absolute Time to Geometric Unity},
	GitHub Repository, 2025.
	[Online]. Available: \href{https://github.com/jpascher/T0-Time-Mass-Duality}{https://github.com/jpascher/T0-Time-Mass-Duality}.
	
\end{thebibliography}